\section{Related Work}

\subsection{Dynamic Tree Visualization}

A substantial body of work has studied how to visualize hierarchical structures that evolve over time, with particular emphasis on preserving structural readability and supporting users' mental continuity across timesteps. One representative direction focuses on stable space-filling layouts for evolving hierarchies. Temporal Treemaps transforms a sequence of trees with known correspondences into a static nested treemap, allowing users to inspect appearance, disappearance, merge, and split events without relying solely on animation~\cite{Kopp2018TemporalTreemaps}. EvoCells similarly proposes a treemap layout strategy tailored to evolving tree data, explicitly mapping tree updates to rectangle updates in order to improve layout stability under hierarchy changes~\cite{Scheibel2018EvoCells}. Stable Voronoi treemap approaches further seek to reduce unnecessary layout variation when hierarchies evolve, thereby improving cross-time comparability while preserving the overall spatial organization of the hierarchy~\cite{Hahn2014StableVoronoiTreemaps}.

Another line of work explores static or semi-static representations that directly encode hierarchy changes across multiple steps. Dynamic indented plots visualize changes between successive hierarchies using time-to-space mappings and explicit visual cues for structural updates, making it possible to compare multiple timesteps in a single view~\cite{Burch2014DynamicIndentedPlots}. TreeVersity2 constructs dynamic hierarchies and supports change exploration through space-filling visualizations and multilevel summaries, enabling users to inspect created, removed, and modified nodes over time~\cite{GuerraGomez2013TreeVersity2}. Beyond tree-specific visualizations, Vehlow et al.\ proposed a matrix-based visualization that combines hierarchical grouping information, temporal ordering, and flow-based transition cues to support the analysis of dynamic hierarchies in evolving graph sequences~\cite{Vehlow2016DynamicHierarchiesGraphSequences}. These methods are effective for presenting temporal continuity and preserving structural readability, but they are primarily designed for showing how hierarchies change across time rather than for explicitly modeling how one tree transforms into the next. In addition, when evolution sequences are long and critical changes are sparse, animation- or snapshot-based views often make it difficult to review distant steps, locate salient transitions, or focus on topology-aware differences while retaining essential structural context.

\subsection{Tree Differencing and Change Modeling}

Another relevant line of research focuses on computing structural differences between trees. The classical tree edit distance formulation introduced by Zhang and Shasha established a foundational model for computing the minimum-cost sequence of insert, delete, and relabel operations required to transform one tree into another~\cite{Zhang1989TreeEditDistance}. This formulation has had lasting influence because it provides a principled way to describe structural differences and change processes between two trees. Later work improved the efficiency and robustness of tree edit distance computation. RTED introduced a robust strategy selection mechanism that performs well across different tree shapes~\cite{Pawlik2011RTED}, while AP-TED further improved efficiency and scalability for practical tree comparison~\cite{Pawlik2015APTED}. Approximate or alternative similarity measures, such as pq-Gram distance, have also been proposed to make tree comparison more scalable for large datasets~\cite{Augsten2010PQGram}.

Beyond generic tree edit distance, several methods address differencing in specific classes of structured data. X-Diff detects changes in XML documents by treating them as unordered trees and generating change descriptions based on structural matching~\cite{Wang2003XDiff}. Earlier work on change detection in hierarchically structured information argued that meaningful change descriptions may require richer structural operations and explicit matching strategies~\cite{Chawathe1996ChangeDetection}. Tree differencing has also been extensively studied in source code analysis, where abstract syntax trees provide a natural hierarchical representation of programs. ChangeDistiller extracts fine-grained structural edits from ASTs and organizes them into semantically meaningful change types~\cite{Fluri2007ChangeDistiller}. GumTree further improves AST differencing by combining top-down and bottom-up matching strategies to recover node correspondences and edit scripts that better reflect developer intent, including move operations~\cite{Falleri2014GumTree}. More specialized topology-driven differencing methods, such as branch-based edit distances for merge trees, show how structural comparison can be adapted to topology-centric settings in scientific visualization~\cite{Bauer2022MergeTreeEditDistance}.

Although this literature provides an essential computational foundation, it does not by itself solve the visual analytics problem addressed in our work. Most tree differencing methods are designed to compute distances, correspondences, or edit scripts, rather than to support human-centered analysis of tree evolution. They are therefore less concerned with how differences should be presented to analysts, how unchanged context should be preserved to maintain interpretability, or how local structural edits should be integrated into the analysis of long evolution sequences. These methods are powerful for modeling structural change, but they are not sufficient for topology-aware visual comparison and multi-step evolution analysis.

\subsection{Temporal Summarization and Sequence Abstraction}

Our work is also related to visual analytics research on temporal sequence summarization. A large body of work in this area studies how to abstract long or complex event sequences into more compact representations that reveal high-level patterns, salient intervals, and representative stages. LifeFlow introduced an overview representation for event sequences that aggregates many temporal records into aligned visual timelines, making it easier to identify common pathways, temporal spacing, and outcome-related patterns~\cite{Wongsuphasawat2011LifeFlow}. EventFlow extended this direction by providing simplification operations and a visual complexity model, enabling analysts to progressively reduce clutter and discover important temporal structures in large collections of sequences~\cite{Monroe2013EventFlow}. Outflow uses flow-based visual representations to summarize temporal pathways and transitions across many event sequences~\cite{Wang2014Outflow}, while DecisionFlow addresses high-dimensional temporal event data through milestone-based abstractions and interactive multi-view exploration~\cite{Xu2014DecisionFlow}.

Subsequent work further explored multilevel abstraction, stage analysis, and progression summarization for complex temporal data. EventThread and ET2 focus on stage analysis, temporal segmentation, and progression understanding, providing methods to identify latent stages and summarize how sequences evolve across them~\cite{Zhao2018EventThread,Zhao2018ET2}. Sequen-C supports multilevel overviews of temporal event sequences, allowing analysts to move between coarse and fine-grained representations while preserving sequence-level context~\cite{Wang2022SequenC}. Other systems, such as SSRDVis and LoLo, emphasize pattern summarization, anomaly detection, and multilevel exploration for long event sequences with many events~\cite{Yang2019SSRDVis,Isenberg2025LoLo}. This literature is highly relevant to Tree Evolution because long evolution processes often contain only a limited number of critical structural changes, while the rest of the sequence mainly provides contextual continuity.

However, most temporal summarization methods operate on event records, categorical timelines, or abstract sequence tokens rather than on hierarchical structures with rich topology. As a result, they typically abstract away the hierarchical organization and fine-grained structural semantics that are central to Tree Evolution analysis. Even when they reveal progression stages effectively, they do not naturally support topology-aware comparison between adjacent trees, nor do they preserve the structural context needed to interpret local changes in a hierarchy. Therefore, while temporal sequence summarization offers important ideas for long-sequence abstraction, additional mechanisms are needed to integrate such abstraction with structural change modeling and visual comparison in evolving trees.

\subsection{Domain-Specific Analysis of Evolving Trees}

Several studies have demonstrated the practical importance of analyzing evolving tree-like structures in specific domains. In database systems, Picasso visualizes optimizer behavior across query parameter spaces and helps users inspect how query plans vary under different conditions~\cite{Sudarshan2010Picasso}. More recently, QOVIS directly supports understanding and diagnosing query optimization by exposing intermediate query plans, transformation logic, and the optimization process through a visualization-assisted framework~\cite{Tian2025QOVIS}. These systems are particularly relevant to our work because query plans are a representative example of evolving tree structures, and because the analysis tasks they support---such as understanding transformation behavior, locating problematic steps, and diagnosing abnormal outcomes---closely match the goals of Tree Evolution analysis.

In software engineering and program analysis, hierarchical program representations also evolve over time. Tree-based version control in Envision explores how AST-aware differencing and merging can support semantically meaningful version management~\cite{Schneider2015EnvisionVC}. Related work on dynamic call graphs and call graph evolution further illustrates how tree-like or hierarchical dependency structures can change across program versions or executions, and how visual or analytical methods can help identify recurring patterns, unstable components, or noteworthy structural changes~\cite{Praun2012DynamicCallGraphs,Bao2022CallGraphEvolution}. Tree evolution is also important in search, constraint solving, and scientific data analysis. Visual Search Tree Profiling provides tools for inspecting large search trees, comparing search behaviors, and identifying critical or repeated substructures during search processes~\cite{Dubois2015SearchTreeProfiling}. In scientific visualization, Temporal Merge Tree Maps summarizes the evolution of topological features over time by aligning and linearizing merge trees extracted from temporal scalar data~\cite{Kopp2022TemporalMergeTreeMaps}.

These domain-specific systems provide strong evidence that evolving hierarchical structures are analytically important across many applications. At the same time, they also highlight an important limitation of prior work: their reliance on application-specific assumptions. Query optimizer analysis depends on optimizer rules and plan semantics; AST and call graph analysis depends on programming-language structure and software engineering tasks; search tree profiling depends on domain-specific notions of search behavior and performance; topology-based scientific visualizations depend on properties of scalar fields and feature tracking. These assumptions make such systems highly effective in their target domains, but also limit their generalizability. In contrast, our goal is to develop a general topology-aware visual analytics framework for Tree Evolution that supports corresponding change understanding, context-preserving structural comparison, and long-sequence summarization across different types of tree-structured processes.

\stitle{Summary}
In summary, prior work has contributed important foundations from the perspectives of dynamic hierarchy visualization, tree differencing, temporal sequence abstraction, and domain-specific evolution analysis. Dynamic visualization methods help present changing hierarchies over time, but usually emphasize continuity and readability over explicit transformation understanding. Tree differencing methods provide computational tools for recovering correspondences and edit relationships, but are not designed to support human-centered, context-preserving comparison in long evolution sequences. Temporal summarization methods enable abstraction of long processes, but typically suppress the hierarchical topology and fine-grained structural changes that matter in Tree Evolution analysis. Domain-specific systems demonstrate the real-world importance of evolving trees, but rely on task-specific assumptions that hinder generalization. Our work addresses these limitations by integrating corresponding change modeling, difference-aware substructure extraction, and sequence-level simplification into a unified topology-aware visual analytics framework for Tree Evolution.